\section{Sự cố gặp phải và cách xử lý}
\label{sec:loi}

\subsection{Không truy cập được HTTPS do Security Group thiếu cổng 443}

\textbf{\large Triệu chứng:} Sau khi bật HTTPS (Phần F), mở
\cmd{http://demo.hcmus.edu.vn} trên trình duyệt ở máy cá nhân thì bị chuyển
sang \cmd{https://} rồi \textbf{treo cho đến khi hết thời gian chờ} (connection
timed out), không hiện trang. Trong khi đó, trên chính EC2,
\cmd{curl -kI https://demo.hcmus.edu.vn/} vẫn trả \cmd{HTTP/1.1 200 OK}.

\textbf{\large Phân tích, khoanh vùng:} Kiểm tra lần lượt từ trong ra ngoài theo
đường đi của gói tin (Hình~\ref{fig:arch}):
\begin{enumerate}
  \item \textbf{Apache:} \cmd{apache2ctl configtest} báo \texttt{Syntax OK};
        \cmd{ss -tlnp | grep :443} cho thấy \cmd{apache2} đang LISTEN trên
        \cmd{*:443}; \cmd{curl} nội bộ trả 200 $\Rightarrow$ Apache và SSL hoạt động đúng.
  \item \textbf{Tường lửa trong OS:} \cmd{ufw status} có rule
        \texttt{Apache Full} (= 80 + 443) $\Rightarrow$ ufw không chặn.
  \item \textbf{Nhật ký Apache:} \cmd{tail -f /var/log/apache2/demo\_access.log}
        trong lúc tải lại trang: chỉ ghi nhận request tới cổng 80 (bị
        redirect 301), \textbf{không có} request HTTPS nào từ máy cá nhân tới
        được Apache; \cmd{demo\_error.log} cũng không có lỗi SSL
        $\Rightarrow$ gói tin bị chặn \emph{trước khi} tới máy ảo.
  \item \textbf{Kiểu lỗi:} \emph{timeout} (gói tin bị bỏ đi âm thầm) chứ không
        phải \emph{connection refused} (máy trả RST vì không có tiến trình
        nghe). Đây là dấu hiệu điển hình của một tường lửa ở giữa.
\end{enumerate}

\textbf{\large Nguyên nhân:} Trên AWS, lưu lượng vào EC2 bị lọc bởi \textbf{hai
lớp độc lập}: Security Group (ở hạ tầng AWS) và ufw (trong Ubuntu). Security
Group lúc tạo instance chỉ có inbound cho cổng 22 và 80. Mở cổng 443 trên ufw
là chưa đủ vì gói tin HTTPS đã bị Security Group loại bỏ trước khi đến được
card mạng của máy ảo. Cổng 80 vẫn mở nên trình duyệt nhận được 301, nhưng bước
tiếp theo sang 443 thì bị chặn, khiến lỗi trông như ``trang HTTPS bị treo''.

\textbf{\large Cách sửa:} Vào \textit{EC2 $\to$ Instances $\to$ Security $\to$
Security groups (launch-wizard-9) $\to$ Edit inbound rules}, thêm rule
\textbf{HTTPS, TCP 443, nguồn 0.0.0.0/0}. Không cần khởi động lại Apache hay máy
ảo; tải lại trang thì truy cập được ngay (trình duyệt cảnh báo chứng chỉ tự ký,
đây là điều bình thường trong lab). Kết quả ở Hình~\ref{fig:sg}.

\textbf{\large Bài học:}
\begin{itemize}
  \item Trên cloud, ``mở firewall'' nghĩa là mở \textbf{cả hai lớp}: firewall
        trong OS và firewall của nhà cung cấp (Security Group).
  \item \cmd{curl} trên localhost chạy được không chứng minh được dịch vụ truy
        cập được từ ngoài; cần kiểm thử từ máy khác.
  \item Phân biệt \emph{timeout} và \emph{connection refused} giúp khoanh vùng
        nhanh: timeout $\to$ firewall chặn; refused $\to$ dịch vụ không lắng nghe.
  \item Access log không có dòng nào cho request là bằng chứng rõ ràng request
        không tới được Apache.
\end{itemize}

\subsection{Các điểm phát hiện thêm khi rà soát cấu hình}

\textbf{\large Dư dấu \texttt{/} khi proxy:} JSON trả về ở Hình~\ref{fig:api} có
\texttt{"path": "//hello"}. Nguyên nhân là
\cmd{ProxyPass /api http://127.0.0.1:3000/}: vế trái không có \texttt{/} cuối
còn vế phải có, nên \cmd{/api/hello} được ghép thành
\cmd{http://127.0.0.1:3000//hello}. Cách sửa là để hai vế đồng nhất:
\begin{minted}{apacheconf}
ProxyPass        /api/ http://127.0.0.1:3000/
ProxyPassReverse /api/ http://127.0.0.1:3000/
ProxyPass        /app/ balancer://appcluster/
ProxyPassReverse /app/ balancer://appcluster/
\end{minted}

\textbf{\large Lỗi chính tả trong unit systemd:} Bản \cmd{nodeapi@.service} lúc
làm lab ghi \cmd{WantedBy=mutil-user.target} thay vì \cmd{multi-user.target}.
Các service vẫn chạy vì được khởi động thủ công bằng
\cmd{systemctl enable --now}, nhưng target sai tên nên khi EC2 khởi động lại,
backend sẽ không tự chạy và \cmd{/api}, \cmd{/app} trả 503. Cách sửa: đổi lại
thành \cmd{multi-user.target} (bản nộp kèm ở File~\ref{lst:service} đã sửa),
chạy \cmd{systemctl daemon-reload} rồi
\cmd{systemctl reenable nodeapi@300\{0,1,2,3\}}.
